\section{Architecture}
\label{sec:architecture}

\subsection{Stack}

\begin{table}[H]
\centering
\small
\begin{tabular}{@{}lll@{}}
\toprule
Component & Version & Role \\
\midrule
\code{cosmos/evm} & v0.7.3 & EVM execution, JSON-RPC, fee market, precompiles \\
\code{cosmos/cosmos-sdk} & v0.54.2+ & accounts, bank, staking, distribution, gov, slashing, upgrades \\
\code{cometbft/cometbft} & v0.39.3 & consensus, p2p, mempool, state sync \\
\code{cosmos/ibc-go} & v11.x & IBC clients, connections, channels, ICS-20 \\
go-ethereum & v1.17 (Cosmos fork) & EVM interpreter, pulled in by \code{cosmos/evm} \\
Go toolchain & 1.26 & build \\
\bottomrule
\end{tabular}
\caption{Target stack at the time of this draft (ENGINEERING \S3). Versions
move with upstream security releases; the policy in \S\ref{sec:neverfork} is
what stays fixed.}
\end{table}

CometBFT provides single-slot finality: a block is final when it is committed,
with no reorganisation and no probabilistic confirmation window. Target block
time is on the order of 1--2 seconds; the tokenomics record currently assumes
$\approx$1.5\,s (21\,038\,400 blocks per year) and will be re-derived from
observed testnet block time before mainnet (TOKENOMICS \S2).

\subsection{Imported versus implemented}
\label{sec:imported}

The central architectural decision is what \emph{not} to write. Everything in
Table~\ref{tab:imported} is consumed from upstream as a Go module dependency and
is neither reimplemented nor modified. The Konstellation repository itself
contains only the items in Table~\ref{tab:implemented} (ENGINEERING \S7).

\begin{table}[H]
\centering
\small
\begin{tabularx}{\textwidth}{@{}LL@{}}
\toprule
Provided upstream & Package \\
\midrule
EVM execution and state database & \code{cosmos/evm} \code{x/vm} \\
ERC-20 $\leftrightarrow$ IBC token representation & \code{cosmos/evm} \code{x/erc20} \\
EIP-1559 fee market & \code{cosmos/evm} \code{x/feemarket} \\
Precompiles: staking, distribution, bank, gov, ICS-20, werc20, bech32, p256 & \code{cosmos/evm} \code{precompiles/} \\
Ethereum JSON-RPC and WebSocket & \code{cosmos/evm} \\
EIP-712 signing (MetaMask signs Cosmos messages) & \code{cosmos/evm} \\
Staking, distribution, governance, slashing, bank, auth, feegrant, authz & \code{cosmos-sdk} \\
Upgrade orchestration, evidence & \code{cosmos-sdk} \\
Circuit breaker & \code{cosmos-sdk} \code{contrib/x/circuit} (\dref{D14}) \\
Consensus, p2p, mempool, state sync, snapshots & \code{cometbft} \\
IBC clients, connections, channels, ICS-20 & \code{ibc-go} \\
\bottomrule
\end{tabularx}
\caption{Imported components. None is forked.}
\label{tab:imported}
\end{table}

\begin{table}[H]
\centering
\small
\begin{tabularx}{\textwidth}{@{}lL@{}}
\toprule
Implemented in \code{konstellation} & Notes \\
\midrule
\code{app.go} wiring and ante-handler chain & the order of ante decorators is security-critical and is in audit scope \\
\code{x/compliance} & chain-wide allow/block lists, timelock, emergency freeze, precompile, IBC receive gate (\dref{D6}, \S\ref{sec:compliance}) \\
\code{x/ratelimit} & IBC value-flow quotas; written because no rate limiter exists for \code{ibc-go} v11 (\dref{D15}, \S\ref{sec:rails}) \\
Issuance function & the \dref{D4} curve, installed as the stock \code{x/mint} module's \code{MintFn} hook --- not a module \\
Base-fee burn & an application-level EndBlock step (\dref{D5}) \\
Genesis parameter set & every parameter in \S\ref{sec:tokenomics} \\
Chain identity & EIP-155 IDs, Cosmos chain-id, Bech32 prefix, denomination \\
Upgrade handlers & one package per release, kept permanently \\
Genesis preinstalls & the contracts in \S\ref{sec:aa}, verified against pinned bytecode at \code{init} \\
\bottomrule
\end{tabularx}
\caption{Custom code. A near-empty \code{x/} directory is the success signal:
every file there is code that must be audited once and patched forever.}
\label{tab:implemented}
\end{table}

\subsection{The never-fork policy}
\label{sec:neverfork}

\code{cosmos/evm}, \code{cosmos-sdk}, \code{cometbft} and \code{ibc-go} are
consumed as Go module dependencies only. There is no Konstellation fork of any
of them, no repository by those names in the organisation, and no local
\code{replace} directive in the build pointing at anything the organisation
controls. Dependencies are pinned to release tags and commit hashes, never to
branch names, so that an upstream force-push cannot silently change what
validators execute (ENGINEERING \S2.1--\S2.3).

The reason is operational, and it is drawn from a specific incident. In August
2026 a critical balance-underflow bug in \code{cosmos/evm}
(GHSA-7g4w-cg88-2cq2) drained six Cosmos EVM chains of roughly USD~5.72\,M.
The fix had existed on the upstream \code{main} branch since 15~May; it was not
backported to release branches until 19~August. Chains running forks had to
notice the patch, port it across their fork, build, test and coordinate a
state-breaking upgrade. One affected network's post-mortem stated that twenty
hours was not a realistic window to coordinate such an upgrade across
thirty-eight validators. Earlier that year a nested-execution bug
(ASA-2026-002, CVSS~9.3) had cost another network approximately USD~7\,M.
Cosmos Labs later disclosed that it had shipped fixes for thirty-seven
vulnerabilities over thirteen months without publicly describing the exploit
paths, and two of the relevant release notes described the security backport
only as ``important security fixes'' (ENGINEERING \S2.1, \S4).

The conclusion is that every fork hop between upstream and a running chain adds
latency to security patches, and that latency is what the exploits used. The
policy that follows:

\begin{itemize}
\item Every upstream patch release is treated as a security release until
  proven otherwise, and its diff is reviewed by hand with attention to the EVM
  state database and precompiles.
\item The engineering target is to build, test and be ready to ship within
  \textbf{24 hours} of any upstream tag. An automated watcher opens a tracked
  issue within six hours of a new upstream release; an unassigned issue after
  one working day is treated as an incident (ENGINEERING \S4.2, \S17).
\item \code{govulncheck} runs on every change and nightly, because a new
  advisory can land against code nobody touched; every accepted finding is
  justified in writing and re-reviewed on every upstream bump (ENGINEERING \S4.3).
\item Custom behaviour is added by composition --- registering precompiles into
  the precompile map, composing the ante chain, wrapping the mempool --- never by
  patching upstream. Where upstream lacks something (the IBC rate limiter for
  \code{ibc-go} v11) it is written as a separate module against the upstream
  interfaces (\dref{D15}).
\end{itemize}

A small validator set is, in this one respect, a security advantage: a
coordinated binary swap across ten nodes takes under an hour.

\subsection{Execution: sequential at launch}
\label{sec:blockstm}

\code{cosmos/evm} v0.7 introduced BlockSTM, parallel transaction execution
under a software-transactional-memory scheduler, bundled with a new
``virtual fee collection'' path. Both are opt-in. Konstellation launches with
\textbf{both switched off} and uses the sequential runner (ENGINEERING \S2.5,
\S7.3). The rationale:

\begin{itemize}
\item Both critical 2026 bugs were balance-accounting bugs in the
  \emph{sequential} path, which has years of production mileage. The parallel
  path has months.
\item A new chain does single-digit transactions per second; parallel execution
  matters in the hundreds.
\item Backing parallel execution out after launch, with funds on chain, would be
  an emergency state-breaking downgrade. Enabling it later is a planned
  governance upgrade with a tested rollback.
\end{itemize}

The enablement plan is to run a non-validating \emph{shadow node} with the
parallel runner from day one --- same network, same genesis, zero voting power ---
and compare its application hash against the validators' every block. A
divergence finds a BlockSTM bug with no funds at risk. After a clean month,
validators are switched by governance upgrade (launch-sequence phase~10,
\S\ref{sec:launch}).

The application-side EVM mempool (``Krakatoa'') is independent of BlockSTM and
is \textbf{on}, matching the upstream default (\dref{D13}).

\subsection{Node topology}

Validators sit behind sentry nodes and have no public IP address; peer exchange
is disabled on validators and the validator's node identity is never gossiped
by its sentries. Validator signing keys are held by threshold signers
(Horcrux across three or more regions is the preferred configuration for a chain
holding user funds), which both provides high availability and makes
double-signing cryptographically hard. Release binaries are never built on a
validator: they are produced by continuous integration from a signed tag, built
twice independently, and published only if both builds agree on the checksum
(ENGINEERING \S2.6, \S9). Operators are diversified by provider, autonomous
system and jurisdiction, mixing bare-metal and cloud hosts.
